\subsection{The complete twelve-arrow triangular exclusion}

\begin{lemma}[A singleton return hom-set]\label{C:singleton-return}
Let $C$ be a category with two objects $a,b$, with
$\Hom(b,a)=\{t\}$ and $\Hom(a,b)\ne\varnothing$.
Then $BC$ is contractible. More generally, the same conclusion holds
for any category with an object $a$ such that every
$\Hom(b,a)$, $b\ne a$, is a singleton and at least one
$\Hom(a,b)$, $b\ne a$, is nonempty. The opposite criterion also holds.
\end{lemma}
\begin{proof}
Put $M=\End(a)$ and choose $s:a\to b$, where $b\ne a$.
If $t:b\to a$ is the unique return arrow, every $m\in M$
satisfies $mt=t$. Thus $e=ts$ satisfies $me=e$ for all $m$.
It is the component of a natural transformation from the constant
functor to the identity functor on $M$, so $BM$ is contractible.
For the full inclusion $i:\{a\}\to C$, the comma category
$(a\downarrow i)$ has initial object $1_a$. For $b\ne a$,
$(b\downarrow i)$ has its single object $t$ and endomorphism
monoid $M$, since $mt=t$. Every comma category is therefore
contractible. Quillen's Theorem A gives $BM\simeq BC$.
Applying the same proof to $C^{\mathrm{op}}$ gives the dual assertion.
No finiteness assumption is needed.
\end{proof}

\begin{remark}[Retract pruning in a twelve-arrow search]\label{C:twelve-retract-pruning}
If $u:x\to y$ and $v:y\to x$, $x\ne y$, satisfy $vu=1_x$,
the full subcategory obtained by deleting $x$ is Cauchy dense.
Adding or removing a retract preserves the classifying space, as in
Lemma~\ref{C:two-object-cup-pruning}. At twelve arrows this smaller
category has at most eleven arrows, so the completed eleven-arrow
exclusion rules out $\mathbb CP^2$. In particular a trivial
endomorphism vertex with arrows both ways can be pruned before any
composition is chosen. There is one further deletion.  Suppose a vertex $x$ has exactly
one arrow from $x$ to any of the other vertices combined, say
$x\to y$, and $\End(y)=\{1_y\}$.  Its own endomorphisms and
incoming arrows are unrestricted.  For the full inclusion deleting
$x$, the comma category $(x\downarrow i)$ is the single arrow
$x\to y$ with endomorphism monoid $\End(y)=\{1_y\}$.
Retained vertices have identity initial objects in their comma
categories.  Quillen's Theorem A therefore permits deletion of $x$.
The criterion with incoming arrows is dual.  The file
\texttt{c\_categories12\_cofinal.cpp} uses precisely this stronger
deletion; the preceding \texttt{reduced} version used only its
special case for a leaf with trivial endomorphisms.
The other reductions are used in the new
\texttt{c\_categories12\_minimal.cpp} search; they do not alter the
recorded scopes or counts of the earlier enumerations.
\end{remark}

\begin{lemma}\label{C:twelve-triangular}
No triangular two-object category with at most twelve arrows has
classifying space homotopy equivalent to $\mathbb CP^2$.
\end{lemma}
\begin{proof}
Let $(m,n,s)$ denote the two endomorphism orders and the nonzero
cross-hom cardinality, with $m\geq n$. The preceding eleven-arrow
proof and Lemma~\ref{C:two-object-cup-pruning} leave, at twelve arrows,
only
\[
 (8,1,3),\ (7,2,3),\ (6,3,3),\ (5,4,3),\
 (5,3,4),\ (5,5,2),
\]
or act-cone cases over at most seven elements. To see that this list
is exhaustive, singleton cross-hom sets are excluded by the join
description. For two cross arrows a vertex of order at most four
is rationally acyclic and the two-point argument applies; the only
remaining endomorphism split is $(5,5)$. For $s\geq3$, two vertex
orders at most four again give zero rational products. One may
therefore assume $m\geq5$, and enumerate $m+n+s=12$ with $m\geq n$.
All cases with $n=1,m\leq7$ are covered for arbitrary $s$ by
Proposition~\ref{C:seven-act-cones}. That proposition also covers
$(7,2,3)$, because each two-element monoid algebra is semisimple
over $\mathbb Q$.

Here is an independent complete enumeration for $(8,1,3)$.
If $H^2(BM;\mathbb Q)=0$, the relative-cup argument already
excludes the act cone. Otherwise the minimal ideal has at least
two rows and two columns. Its local group must be trivial: a
nontrivial group would require at least eight ideal elements and
a further global identity. Its band is consequently $2\times2$,
$2\times3$, or $3\times2$.

For a $2\times2$ band, label its elements $x_{ij}$ and write the
other elements as $1,a,b,c$. Each extra element specifies a row
and column transformation on two points. The balanced presentation
for rational $H_2$ forces equality of their augmentation characters.
There are exactly six such translation pairs. Choose their types
in nondecreasing order; this removes only permutations of $a,b,c$.
There are $\binom{8}{3}=56$ type profiles. The nine products among
$a,b,c$ may be any of the eight elements with the prescribed
composite translations. Exhaustive associativity propagation on
these choices leaves 4032 associative tables and 492 ordinary
isomorphism classes. Isomorphism recognition uses exactly the
permutations of the two rows, two columns, and three extra elements;
the minimal ideal is characteristic, so this is the complete
isomorphism test in the chosen coordinates.

For a $2\times3$ band there is one extra element. There are
$4\cdot27=108$ translation pairs, and its square is tried at
every compatible table entry. This gives 56 associative tables
and twelve ordinary isomorphism classes. The opposite orientation
gives another twelve classes. Thus the complete list has 516
tables. The program \path{code/c_eight_H2_targets.cpp} constructs
this list without using a published census. The independent Python
check verifies every identity and all $8^3$ associativity equations
of each table. Its exact rational balanced-module ranks discard
six tables with $H_2(BM;\mathbb Q)=0$, leaving 510 targets.

For each of the 510 targets, enumerate all right actions on three
points. Choose a generating set, assign every transformation to
each generator subject only to its exact power relations, extend
along generator words, and verify all multiplication equations.
State permutations remove precisely isomorphic actions. Because
ordinary opposite tables are included, both arrow orientations
are covered. There are 8828 action classes. The exact normalized
category computation excludes 2266 by $b_1(-;\mathbb F_2)$ and
3321 by $b_2(-;\mathbb F_2)$. Of the others, 3226 have an explicit
degree-three primitive for the degree-two generator's square.
The remaining fifteen have third mod-two Betti number one, and
therefore are not $\mathbb CP^2$. The tables, actions, cocycles,
primitives, and complete residual ranks are retained in
\path{code/outputs/c_eight_all_act_cones.json}.

It remains to treat the other four profiles in the displayed list.
For $(6,3,3)$, the relative-cup argument again restricts the larger
monoid to those with nonzero rational $H^2$. Its minimal band is
$2\times2$ with one extra element. The same translation-pair and
square enumeration gives exactly five ordinary isomorphism classes.
For $(5,4,3)$ and $(5,3,4)$, the larger monoid must be the
five-element sphere monoid. In $(5,5,2)$, both monoids must be that
monoid: every other five-element monoid has rationally acyclic
classifying space and is covered by the two-point reduction.

The own three- and four-element monoid censuses, with opposites
restored, supply respectively seven and 35 ordinary isomorphism
classes for the smaller vertex. The program
\path{code/c_twelve_triangular_bisets.py} enumerates every right
action of the larger monoid up to state permutation and \emph{all}
conjugates of every left action of the smaller monoid. It retains
exactly the commuting pairs, checking commutation first on
generators and then on every pair of elements. Thus no relative
position of the two actions is lost. Its complete counts are
\[
\begin{array}{c|r|r|r|r|r}
(m,n,s)&\text{bisets}&b_1\ne0&b_2\ne1&x^2=0&x^2\ne0\\\hline
(6,3,3)&1740&460&890&390&0\\
(5,4,3)&2264&572&1244&448&0\\
(5,3,4)&2102&851&759&491&1\\
(5,5,2)&8&1&1&6&0.
\end{array}
\]
The tests are applied in the indicated order, so the columns are
disjoint and sum to the biset count. The sole nonzero-square
residual has third mod-two Betti number one. Every asserted zero
square has a full cochain primitive checked on all normalized
four-simplices. All four profiles are therefore excluded, which
completes the proof.
\end{proof}
